\section{幂级数、递推与收敛真题}
\problem{74}{递推微分方程后求幂级数和}{2009 年第 1 届，非数学类，七，\src{01}}
$u_n'=u_n+x^{n-1}\ee^x$，$u_n(1)=\ee/n$。求 $\sum_{n=1}^\infty u_n(x)$ 及定义区间。
\sol 乘积分因子 $\ee^{-x}$，得 $(\ee^{-x}u_n)'=x^{n-1}$；用初值积分得 $u_n=\ee^x x^n/n$。因此
\[
\sum_{n=1}^\infty u_n(x)=-\ee^x\ln(1-x),\quad |x|<1.
\]
在 $x=-1$ 交错收敛，级数和由 Abel 定理得 $-\ee^{-1}\ln2$；$x=1$ 项为 $\ee/n$ 发散；$x<-1$ 通项不趋零。因此收敛集合 $\ans{[-1,1)}$。

\problem{75}{从几何级数求有理函数和}{2011 年第 3 届，非数学类，一(4)，\src{03}}
求 $S(x)=\sum_{n=1}^\infty(2n-1)x^{2n-2}/2^n$ 的和与定义域，并求 $\sum_{n=1}^\infty(2n-1)/2^{2n-1}$。
\sol 令 $q=x^2/2$，用 $\sum q^{n-1}=(1-q)^{-1}$ 及其导数，得
\[
S(x)=\frac12\sum_{n=1}^\infty(2n-1)q^{n-1}
=\frac{1+q}{2(1-q)^2}=\frac{2+x^2}{(2-x^2)^2}.
\]
根值检验得 $|x|<\sqrt2$，端点通项不趋 0。所求数值为 $S(1/\sqrt2)=\ans{10/9}$。

\problem{76}{二阶导数控制函数值级数}{2013 年第 5 届，非数学类，二，\src{05}}
$f$ 在 0 二阶可导，且 $\lim_{x\to0}f(x)/x=0$。证明 $\sum_{n=1}^\infty|f(1/n)|$ 收敛。
\sol 极限先给出 $f(0)=0$。由差商定义 $f'(0)=0$。二阶 Taylor--Peano 式给
$f(x)=\frac12f''(0)x^2+o(x^2)=O(x^2)$，所以 $|f(1/n)|\le C/n^2$ 对充分大的 $n$ 成立。比较 $p$ 级数，结论成立。这里二阶可导足以有 Peano 展开，不需要假设 $f''$ 连续。

\problem{77}{调和数与非负双重求和}{2013 年第 5 届，非数学类，七，\src{05}}
求 $\displaystyle\sum_{n=1}^\infty\frac{H_n}{(n+1)(n+2)}$，其中 $H_n=\sum_{k=1}^n1/k$。
\sol 项均非负，可先按有限矩形求和，再由单调极限交换次序：
\[
\sum_{n=1}^\infty\sum_{k=1}^n\frac1{k(n+1)(n+2)}
=\sum_{k=1}^\infty\frac1k\sum_{n=k}^\infty\left(\frac1{n+1}-\frac1{n+2}\right)
=\sum_{k=1}^\infty\frac1{k(k+1)}=\ans1.
\]
非负项交换不依赖先验收敛，若换成带符号项则须另证绝对/一致收敛。

\problem{78}{跳跃点傅里叶和}{2015 年第 7 届，非数学类，一(4)，\src{07}}
将 $f(x)=3$（$-5\le x<0$）、$f(x)=0$（$0\le x<5$）作周期 10 延拓，求其傅里叶级数在 $x=0$ 的和。
\sol 周期延拓在 0 的左极限为 3、右极限为 0。Dirichlet 收敛定理给跳跃平均值 $\ans{3/2}$。级数的和并非函数在跳点所定义的 $f(0)=0$；边界约定不改变跳跃平均。

\problem{79}{阶乘分母的整系数拆分}{2015 年第 7 届，非数学类，四，\src{07}}
求 $\displaystyle\sum_{n=0}^\infty\frac{n^3+2}{(n+1)!}(x-1)^n$ 及收敛半径。
\sol 令 $t=x-1,m=n+1$。因为 $(n^3+2)=(m-1)^3+2=m^3-3m^2+3m+1$，且
$\sum m^3t^m/m!=(t^3+3t^2+t)\ee^t$、$\sum m^2t^m/m!=(t^2+t)\ee^t$、$\sum mt^m/m!=t\ee^t$，整理后
\[
\ans{\sum_{n=0}^\infty\frac{n^3+2}{(n+1)!}(x-1)^n
=\frac{(t^3+t+1)\ee^t-1}{t}}.
\]
$t=0$ 取连续延拓，值为 2。阶乘使半径 $\ans{R=\infty}$。可用首两项复核封闭式：常数项 2，$t$ 项 $3/2$。

\problem{80}{对幂级数求导再积分}{2023 年第 15 届，非数学 B 类，四，\src{15B}}
求 $S(x)=\displaystyle\sum_{n=1}^\infty\frac{(-1)^{n-1}x^{2n}}{n(2n-1)}$ 的收敛域与和。
\sol 半径为 1。对 $|x|<1$ 求导两次：
\[
S''(x)=2\sum_{n=1}^\infty(-1)^{n-1}x^{2n-2}=\frac2{1+x^2}.
\]
$S(0)=S'(0)=0$，积分得 $S'(x)=2\arctan x$、$S(x)=2x\arctan x-\ln(1+x^2)$。在端点原级数绝对收敛，因为通项绝对值为 $O(n^{-2})$；端点由连续延拓，故 $\ans{x\in[-1,1]}$。

\problem{81}{部分分式决定 Taylor 半径与端点}{2024 年第 16 届，非数学 B 类，三，\src{16B}}
求 $f(x)=1/[(x-1)(x+3)]$ 在 $x=2$ 的 Taylor 级数及收敛区间。
\sol 令 $h=x-2$，
\[
f(x)=\frac14\left(\frac1{1+h}-\frac1{5+h}\right)
=\frac14\sum_{n=0}^\infty(-1)^n(1-5^{-n-1})h^n.
\]
最近奇点距中心 1，故半径为 1，区间 $(1,3)$。在 $h=1$ 通项不趋零；在 $h=-1$ 原函数有极点且通项不趋零，两个端点都发散。

\problem{82}{递推差分望远镜成正项级数和}{2023 年第 15 届，非数学 A 类，六，\src{15A}}
$x_0=1/3$，$x_{n+1}=x_n^2/(1-x_n+x_n^2)$。求 $\sum_{n=0}^\infty x_n$。
\sol 对 $0<x<1$，有 $0<x^2/(1-x+x^2)<x$，故 $x_n\downarrow0$。直接通分可得
\[
\frac{x_n}{1-x_n}-\frac{x_{n+1}}{1-x_{n+1}}=x_n.
\]
求和望远镜：$\sum_{n=0}^N x_n=\frac{x_0}{1-x_0}-\frac{x_{N+1}}{1-x_{N+1}}\to\ans{1/2}$。

\problem{83}{两个不可公度周期迫使 Fourier 系数消失}{2016 年第 8 届，非数学类，六，\src{08}}
$f$ 可微，且 $f(x+2)=f(x+\sqrt3)=f(x)$。用 Fourier 方法证明 $f$ 为常数。
\sol 将 $f$ 看作周期 2 的连续函数，系数
$c_n=\frac12\int_{-1}^1f(x)\ee^{-in\pi x}\dd x$。由周期 2 的函数平移不变性，对 $\sqrt3$ 作换元可得
$c_n=\ee^{in\pi\sqrt3}c_n$。若 $n\ne0$，$n\sqrt3$ 非整数，故相位不为 1，$c_n=0$。所有非零频率系数消失，傅里叶唯一性（Fejér 定理）推出 $f=c_0$。

\problem{84}{积分定义数列的敛散、渐近与精确和}{2020 年第 12 届，非数学类，七，\src{12}}
$u_n=\int_0^1(1+t^4)^{-n}\dd t$。讨论 $\sum(-1)^{n-1}u_n$、$\sum u_n$、$\sum u_n/n^p$，并求 $p=1$ 时的和。
\sol 被积函数随 $n$ 严格下降，且 $u_n\to0$（由控制收敛，$t>0$ 时点态趋零），所以交错级数收敛。对 $0\le t\le n^{-1/4}$，$(1+t^4)^n\le(1+1/n)^n<\ee$，故 $u_n\ge\ee^{-1}n^{-1/4}$；于是 $\sum u_n$ 发散。
对 $p=1$，非负项 Tonelli 允许求和与积分交换：
\[
\sum_{n=1}^\infty\frac{u_n}n
=\int_0^1\sum_{n=1}^\infty\frac{(1+t^4)^{-n}}n\dd t
=\int_0^1\ln(1+t^{-4})\dd t
=\ln2+4u_1.
\]
最后一步将对数写成 $\ln(1+t^4)-4\ln t$，对第一项分部。又
$u_1=\int_0^1\frac{\dd t}{1+t^4}
=\frac{\sqrt2}{8}[\pi+2\ln(1+\sqrt2)]$，故
$\ans{\sum u_n/n=\ln2+\frac{\sqrt2}{2}[\pi+2\ln(1+\sqrt2)]}$。
对 $p>1$，$0<u_n\le1$ 且 $u_n/n^p\le1/n^p$，故收敛。$p<1$ 时由于 $u_n\ge c n^{-1/4}$，若 $p<3/4$ 级数发散；$3/4\le p<1$ 需要更精确渐近，不能由上述粗下界断言。实际上令 $s=n^{1/4}t$ 并用 Laplace 极限得 $u_n\sim n^{-1/4}\int_0^\infty\ee^{-s^4}\dd s=\Gamma(1/4)/(4n^{1/4})$，故 $\sum u_n/n^p$ 收敛当且仅当 $p>3/4$。在 $p=3/4$ 比较调和级数发散。
